{\centering\Large\bfseries\color{nebrijared}Abstract\par}
\fontsize{11}{13}\selectfont
\vspace{\baselineskip}
This Master's Thesis addresses the fundamental computational bottlenecks inherent to probabilistic astrobiological inference in the era of high-resolution exoplanetary characterization. With current and upcoming spectroscopic observatories---such as the James Webb Space Telescope (JWST) and the Habitable Worlds Observatory (HWO)---generating multivariate, noisy, and highly coupled atmospheric spectra, classical probabilistic graphical models face severe algorithmic ceilings. Specifically, exact inference methods (e.g., Junction Tree, Variable Elimination) suffer from an exponential memory explosion bounded by $\mathcal{O}(N \cdot d^{w+1})$ due to the high treewidth induced by dense photochemical reaction cycles, while stochastic Markov Chain Monte Carlo (MCMC) approximations exhibit slow asymptotic variance convergence constrained to $\mathcal{O}(1/\sqrt{M})$, leading to severe sample divergence when evaluating ultra-rare events such as industrial technosignatures (chlorofluorocarbons and photochemical non-equilibrium anomalies).

\vspace{\baselineskip}

To overcome these barriers, this work designs, formalizes, and evaluates an end-to-end Quantum Bayesian Network (QBN) architecture executed via Quantum Amplitude Estimation (QAE). By mapping planetary conditional probability tables into quantum superposition states across a $2^N$-dimensional Hilbert space using conditional rotation cascades, the model eliminates exponential tree compilation overheads. Furthermore, by exploiting Grover-driven amplitude amplification and Inverse Quantum Fourier Transforms (IQFT), the quantum inference framework achieves a quadratic asymptotic speedup, bounding estimator convergence strictly to $\mathcal{O}(1/M)$. The architecture is systematically validated on a 17-qubit benchmark modeling the habitable-zone sub-Neptune K2-18b, demonstrating complete mathematical fidelity under ideal statevector execution, analyzing thermal relaxation and depolarizing limits under realistic NISQ noise ($T_1, T_2$), and recovering ground-truth posteriors with an 86.44\% error reduction through Mitiq Zero-Noise Extrapolation (ZNE).
